% SPDX-License-Identifier: CC-BY-SA-4.0
% Session: Form parallelism to concurrency
% Exercise: Amdahl's law

\begingroup

\begin{exercice}[Vivacité]
  On cherche à implémenter un registre linéarisable avec une méthode \lstinline{getAndMax}
  qui écrit l'argument s'il est supérieur à la valeur stockée dans le registre, et retourne
  la valeur du registre à l'appel de la méthode.
  \begin{question}
  \item Sous quelles conditions chacun des programmes suivants termine-t-il ? 
  \end{question}

  \begin{minipage}{.47\textwidth}
    \begin{lstlisting}[gobble=6]
      class Maximizer1 {
        AtomicInteger x = new AtomicInteger(0);
        int getAndMax (int v) {
          while(true) {
            int old = x.get();
            if(x.compareAndSet(old, max(old, v)))
            return old;
          }
        }
      }

    \end{lstlisting}
    \begin{lstlisting}[gobble=6]
      class Maximizer3 {
        AtomicInteger x = new AtomicInteger(0);
        int getAndMax (int v) {
          while(true) { 
            int old = x.get() / 2;
            if(v <= old) return old;
            if(x.compareAndSet(2*old, 2*old+1)){
              x.set(2*v);
              return old;
            }
          }
        }
      }
    \end{lstlisting}
  \end{minipage}
  \hspace{\fill}
  \begin{minipage}{.47\textwidth}
    \begin{lstlisting}[gobble=6]
      class Maximizer2 {
        AtomicInteger x = new AtomicInteger(0);
        int getAndMax (int v) {
          while(true) {
            int old = x.get();
            if(v <= old) return old;
            if(x.compareAndSet(old, v))
            return old;
          }
        }
      }
    \end{lstlisting}
    \begin{lstlisting}[gobble=6]
      class Maximizer4 {
        AtomicInteger x = new AtomicInteger(0);
        int getAndMax (int v) {
          while(true) { 
            int old = x.get() / 2;
            if(x.compareAndSet(2*old, 2*old+1)){
              x.set(2*max(old, v));
              return old;
            }
          }
        }
      }

    \end{lstlisting}
  \end{minipage}  
\end{exercice}

\endgroup
\endinput
